\subsection{The scalar-complex exclusion in every ordinary type-A rank}

The preceding elementary proof is independent of the
following extension. A known theorem of Lichtenbaum
supplies its module inequality in every polynomial
dimension. The remaining argument then takes place
directly in the actual ambient ring.

Let \(k\) have characteristic zero, let \(n\ge2\), and
let \(A\) be the standard degree-two polynomial
realization for the ordinary Soergel category of
\(W=S_n\times S_n\). This includes the reflection
realization and the permutation realization with
fixed coordinates. Fix a simple reflection \(s\)
in each factor. For \(\tau=(s,s)\), choose coordinates
\[
 A=k[x,y,z_1,\ldots,z_{N-2}],\qquad
 \tau(x)=-x,\quad \tau(y)=-y,\quad\tau(z_j)=z_j.
\]
Such coordinates exist because the minus eigenspace
has dimension two. Put \(\delta=v+v^{-1}\).

\begin{finding}[\fid{F59}, complete bounded scalar exclusion]
\label{f:allRankScalarClosure}
There is no bounded ordinary Soergel complex \(X\)
with the actual ambient Euler class
\[
 [X]=b_s\otimes b_s+(b_s-\delta)\otimes(b_s-\delta),
\]
and no bounded scalar complex \(C\) of finite graded
\(k\)-vector spaces with \(\chi_v(C)=\delta^2\),
admitting an actual derived bimodule isomorphism
\begin{equation}\label{eq:allRankScalarQuadratic}
 X\otimes_A^{\mathbf L}X\simeq C\otimes_kX.
\end{equation}
Thus the entire bounded scalar-complex alternative,
including nonflat cohomology, is excluded in every
stated ordinary type-A rank.
\end{finding}

\begin{proof}
\par\smallskip\noindent\textbf{The precise known local input.}\quad
For a complete unramified regular local ring \(B\)
and finite modules, Lichtenbaum's
Theorem~2(2), (4), (1), pp.~222--223, respectively gives:
nonnegativity of the higher Tor Euler characteristics
when the relevant lengths are finite;
vanishing of \(\operatorname{Tor}_1\) if
\(\chi_1=0\) and one module is torsion-free;
and vanishing of all later Tor groups after a
positive Tor vanishes \cite{LichtenbaumTor}.
We use it only over an equicharacteristic power
series ring and only with the indicated torsion-free
module. The field-case corollary in
\cite[supplementary lecture April~18, p.~169]{Hochster615}
independently corroborates these hypotheses.
These known results are cited, not asserted as new.

\par\smallskip\noindent\textbf{The all-variable module inequality.}\quad
Let \(A_0=k[x_1,\ldots,x_N]\) be positively graded,
\(\mathfrak m=(x_1,\ldots,x_N)\), let \(V\ne0\) be
a finite-dimensional graded right module, and let
\(M\) be a finite graded left module of rank \(r\).
We prove
\begin{equation}\label{eq:allRankTensorDimension}
 \dim_k(V\otimes_{A_0}M)\ge r\dim_kV,
\end{equation}
with equality exactly when \(M\) is graded free.
The case \(N=0\) is immediate.

First suppose \(M=Q_0\) is torsion-free.
Choose a minimal graded free resolution \(F\).
The \(N\)-variable Koszul resolution of \(k\) gives
\(\operatorname{Tor}_i(k,Q_0)=0\) for \(i>N\).
Minimality identifies this Tor with
\(F_i/\mathfrak mF_i\), so the resolution is finite.
If \(f_i=\operatorname{rank}F_i\), localization
gives \(\sum_i(-1)^if_i=r\).
The positive grading makes the action of
\(\mathfrak m\) on \(V\) nilpotent. Thus every
term of \(V\otimes F\) has finite dimension, and
\begin{equation}\label{eq:allRankChiOne}
 \begin{split}
 \dim(V\otimes Q_0)-r\dim V
 &=\chi_1^{A_0}(V,Q_0),\\
 \chi_1^{A_0}(V,Q_0)
 &=\sum_{i\ge1}(-1)^{i-1}
               \dim_k\operatorname{Tor}_i^{A_0}(V,Q_0).
 \end{split}
\end{equation}

Localize at \(\mathfrak m\) and complete to
\(B=k[[x_1,\ldots,x_N]]\). Completion is flat
\cite[Lemmas~10.97.2--10.97.3, tags 00MB--00MC]{StacksCompletion}.
A finite torsion-free module over the domain \(A_0\)
embeds in a finite free module: choose a fraction-field
basis and clear the denominators of a finite set of
generators. Flatness preserves that injection, so
\(\widehat Q_0=B\otimes_{A_0}Q_0\) is torsion-free.
Some power of \(\mathfrak m\) kills \(V\), so its
extension to \(B\) is canonically the same module.
The completed resolution computes exactly the
same finite tensor complex \(V\otimes F\).
All Tor modules therefore have finite length, equal
to their \(k\)-dimension, and the completed ring
satisfies every hypothesis of the known theorem.
It gives \(\chi_1\ge0\).

If equality holds in \eqref{eq:allRankChiOne},
Lichtenbaum's torsion-free clause gives
\(\operatorname{Tor}_1^B(V,\widehat Q_0)=0\),
and rigidity gives vanishing of every positive Tor.
If \(Q_0\) were nonfree, its finite minimal resolution
would have a last nonzero term \(F_t\), \(t>0\).
The completed final differential still has entries
in the maximal ideal. The nonzero socle of \(V\)
in each generator of \(F_t\) is in its tensor kernel.
There is no incoming term beyond \(F_t\), so
\(\operatorname{Tor}_t^B(V,\widehat Q_0)\ne0\),
a contradiction. Thus equality implies graded
freeness; the converse is immediate.

For arbitrary \(M\), let \(T\) be the full graded
torsion submodule and \(Q_0=M/T\). The quotient is
finite graded torsion-free of rank \(r\).
If it is nonfree, strictness for \(V\otimes Q_0\)
passes through the surjection \(V\otimes M\to
V\otimes Q_0\). If it is free, lift a homogeneous
basis to split \(M\to Q_0\) as a left-module map.
When \(T\ne0\), the remaining tensor \(V\otimes T\)
surjects onto
\[
 (V/V\mathfrak m)\otimes_k(T/\mathfrak mT)\ne0.
\]
Both factors are nonzero by graded Nakayama.
Thus strictness also holds in this case.
Equality holds precisely for a free \(M\).
No bimodule splitting was used.

\par\smallskip\noindent\textbf{Generic cohomology in the full group.}\quad
Return to the actual ambient \(A\).
Write \(M_i=H^i(X)\), \(c_i=\dim_kH^i(C)\),
and \(d(u)=\sum_i c_i u^i\).
Over the ungraded fraction field, exact ordinary
graph localization has actual cohomology character
\[
 F(u)=\sum_{g\in S_n\times S_n}f_g(u)g,
 \qquad f_g(u)\in\mathbb N[u,u^{-1}].
\]
The primary graph model is
\cite[\S\S1.1--1.2, equations (1.1)--(1.2),
Theorem~1.6 and Remark~1.7]{EWLocalized}.
Generic K\"unneth gives \(F^2=dF\), and the
prescribed Euler classes give
\[
                    F(-1)=2(e+\tau),\qquad d(-1)=4.
\]
Augmentation gives \(\sum_gf_g=d\).
For every \(u>0\), the normalized coefficients
form an idempotent probability measure.
Positivity makes its fixed support a finite subgroup,
and the maximal-coefficient weighted-average argument
makes the measure uniform. Equality for all positive
\(u\) is a Laurent-polynomial identity.
The Euler support identifies this subgroup as
\(\{e,\tau\}\). Hence
\begin{equation}\label{eq:allRankUniformGraphs}
 F(u)=\frac{d(u)}2(e+\tau),\qquad
 \operatorname{rank}_A M_i=c_i.
\end{equation}
This takes place in the full group, without a
restriction of the object to a rank-one subcategory.
Internal grading has deliberately been forgotten.

\par\smallskip\noindent\textbf{The actual top module is free.}\quad
For any nonzero finite graded bimodule \(M\), finite
on the left, \(M\otimes_AM\ne0\).
Indeed \(V=k\otimes_AM\) is a nonzero finite graded
right origin module, and the left-origin fiber
of the square surjects onto
\[
 (V/V\mathfrak m)\otimes_k(M/\mathfrak mM)\ne0.
\]
Let \(a=\max H(X)\), \(c=\max H(C)\).
The top-truncation identity and scalar K\"unneth
then force \(c=a\). If \(b\) is the highest generic
degree, its nonzero graph top has nonzero square,
and the generic quadratic forces \(2b=b+c\).
Therefore \(b=c=a\) and \(r=c_a=\operatorname{rank}M_a>0\).
The actual top cohomology isomorphism is
\[
 M_a\otimes_AM_a\simeq H^a(C)\otimes_kM_a.
\]
Taking its left-origin fiber gives equality in
\eqref{eq:allRankTensorDimension}, with scalar
multiplicity \(c_a=r\). Thus \(M_a\) is left-free.

\par\smallskip\noindent\textbf{Graph splitting in the ambient ring.}\quad
For any left-free cohomology module, right multiplication
by \(x,y\) has matrices \(U,V\) satisfying
\[
                         yU=xV,\qquad U^2=x^2I.
\]
These identities hold generically by the two graph
types and globally by torsion-freeness.
Coprimality gives \(U=xZ,V=yZ,Z^2=I\).
Every fixed \(z_j\) acts by its left scalar, again
by generic equality and torsion-freeness.
The bimodule idempotents \((I\pm Z)/2\) have
the two graph actions. Their images are finite
graded projective left modules and are graded free:
lift a homogeneous residue basis to a split free
surjection and use graded Nakayama on its kernel.
They are therefore actual shifted graph modules,
also free on the right.
No spectator variable has been specialized.

\par\smallskip\noindent\textbf{The full descending induction.}\quad
Suppose all \(M_i\), \(i>s\), are free graph modules.
The canonical cycle cross-product is an isomorphism
\begin{equation}\label{eq:allRankNearTopKunneth}
 \bigoplus_{i+j=a+s}M_i\otimes_AM_j
       \longrightarrow H^{a+s}(X\otimes_A^{\mathbf L}X).
\end{equation}
To see this, forget the left action and successively
split off the higher projective cohomology from the
bounded right-projective complex \(X\), cancelling
acyclic top tails and contractible pairs. The remainder
has cohomology at most \(s\), with top \(M_s\).
After tensoring with the second \(X\), the higher
summands contribute \(M_i\otimes M_{a+s-i}\) and
the remainder contributes \(M_s\otimes M_a\) by
top truncation. These are exactly the terms displayed.
Their maps are the canonical cycle cross-product.
The one-sided computation proves its bijectivity;
its intrinsic bimodule linearity makes it a bimodule
isomorphism. No bimodule formality is assumed.

Put
\[
 \phi(N)=\dim_k(N/\mathfrak mN),\qquad
 z=\phi(M_s),\qquad r=c_a>0.
\]
Both \(M_a\otimes M_s\) and \(M_s\otimes M_a\)
have origin-fiber dimension \(rz\).
Graph twists on either side preserve the origin;
all remaining pairs are already free and contribute
\[
 S_s=\sum_{\substack{i+j=a+s\\i>s,\ j>s}}c_i c_j.
\]
Scalar K\"unneth on the right of the actual quadratic
has exceptional terms \(H^a(C)\otimes M_s\) and
\(H^s(C)\otimes M_a\), with dimensions \(rz,rc_s\);
the remaining contribution is the same \(S_s\).
Thus
\[
 2rz+S_s=rz+rc_s+S_s,\qquad
 \phi(M_s)=c_s=\operatorname{rank}M_s.
\]
If \(c_s=0\), graded Nakayama gives \(M_s=0\).
Otherwise its minimal graded free surjection has
rank \(c_s\); its rank-zero kernel inside a free
module over a domain is zero. Thus \(M_s\) is
left-free and, by graph splitting, right-free.
Descending proves that every cohomology module is
a sum of actual shifted identity and \(\tau\) graphs.

\par\smallskip\noindent\textbf{Restored grading and contradiction.}\quad
Only now retain the full global internal grading.
Write the actual cohomology character as
\[
 a(v,u)e+b(v,u)\tau,\qquad
 a,b\in\mathbb N[v,v^{-1},u,u^{-1}],\quad b\ne0,
\]
and put \(d(v,u)=\sum_{i,e}\dim_kH^i(C)_e\,v^eu^i\).
Since every cohomology module is right-projective,
the same one-sided splitting proves canonical
K\"unneth in all degrees. Different graph types
have no maps: their coefficients are killed by
\(2x\). Within a type, the graded origin fiber
determines multiplicities. The actual quadratic gives
\[
                         a^2+b^2=da,\qquad2ab=db.
\]
The Laurent ring is a domain, so \(d=2a\).
The Euler specialization \(u=-1\) gives
\[
                       v^2+2+v^{-2}=2a(v,-1),
\]
contradicting the odd outer coefficients.
Full internal grading has not been retained
illegitimately through generic localization.
\end{proof}

\paragraph{The remaining categorical alternatives.}
There is no residual nonflat-cohomology case under
these bounded ordinary type-A, scalar-coefficient,
ambient-class and actual-quadratic premises.
A single generator already fails its quadratic,
so additional higher-rank relation maps cannot
repair this construction.
Non-scalar central complexes, polynomial coefficient
differentials, changed \(K_0\) embeddings, different
categories such as finite-state kernels or nonproper
sheaf convolution, localizations and unbounded
enlargements remain outside the theorem.
No arithmetic Frobenius comparison, purity, or
designated canonical basis is inferred.

% Separate paragraph for Part III.
\paragraph{Local Tor rigidity and the global Grothendieck dictionary.}
The affine module/sheaf passage is EGA~I,
Theorem~1.3.7, p.~86, Corollary~1.3.8, p.~87,
and Theorem~1.4.1, pp.~90--92.
The bounded free complexes lie in SGA~6, I,
Definition~4.7, pp.~121--122 and Example~4.8, p.~122
\cite{EGA,SGA6}. The equality case of the separately
cited Tor theorem retains positive homological
information; an Euler-zero identity alone would
not imply acyclicity.

The internal grading is a split-torus action in
SGA~3, I, 4.7.2--3, with cote~90, scan~17.
Actual and virtual graded coefficient classes belong
to the representation framework of SGA~6, V,
\S\S2--3, with cote~133, scan~18
\cite{SGA3,SGA6,Archive}.
These are framework anchors, with the existing
source-access qualification for SGA~6, V retained.
They do not furnish the new theorem or a
noncommutative \(\lambda\)-ring.
The ordinary normalization is
\cite[\S\S3.3--3.5]{EW}, and the exact graph
localization was cited in the proof.
The published local theorem is clearly separated
from the new application and descent.
